\section{Statements and background}\label{sec:introduction}

A finite nonempty set $A$ in a Euclidean space is \emph{Ramsey} if, for
every integer $r\ge2$, there exists $N\ge1$ such that every colouring
$\R^N\to\{1,\ldots,r\}$ contains a monochromatic congruent copy of $A$.
The scale is fixed, and the colouring has no regularity assumption.
Ramsey-ness depends only on the distance table, is preserved by nonzero
similarities, and passes to subsets. The assertions concerning circle
sets below apply to circles of positive radius in any Euclidean space:
we first choose an isometric planar representative and then use a
similarity to the unit circle.

Use homogeneous parameters $z=[s:t]\in\Pone(\R)$ for the unit circle:
\begin{equation}\label{eq:circle-map}
 a([s:t])=\left(\frac{s^2-t^2}{s^2+t^2},
                       \frac{2st}{s^2+t^2}\right).
\end{equation}
For nonzero representatives put $[z_i,z_j]=s_it_j-t_is_j$. For four
distinct indices define
\begin{equation}\label{eq:cross-ratio}
 \crat(i,j;k,l)=
 \frac{[z_i,z_k][z_j,z_l]}{[z_i,z_l][z_j,z_k]}.
\end{equation}
All denominators are nonzero. This quantity is unchanged by rescaling
representatives or applying a common element of $\mathrm{PGL}_2(\R)$.
Changing the projective parametrization of the circle therefore does
not change the following field.

\begin{definition}\label{def:moduli}
The \emph{projective moduli field} of a labelled circle set $A$ is
\[
 M(A)=\Q\bigl(\crat(i,j;k,l):i,j,k,l\text{ distinct}\bigr).
\]
A group of at least three of its points is \emph{projectively algebraic}
if all its internal four-point cross ratios are algebraic over $\Q$.
For a group of three points this condition is vacuous.
\end{definition}

Equivalently, a projectively algebraic group is a projective image of a
finite subset of $\Pone(\overline{\Q}\cap\R)$. Indeed, sending three of
its points to $0,1,\infty$ makes every remaining parameter a cross ratio
or a fractional linear transform of one. The same observation shows
that $M(A)$ is generated by the remaining $n-3$ parameters after such a
normalization. In particular, the moduli degree refers to the normalized
parameters, rather than the absolute coordinate field.

\begin{maintheorem}\label{thm:main}
Let $A$ consist of $n\ge7$ distinct points on a circle and suppose that
$\trdeg_\Q M(A)=3$. Then the following are equivalent:
\begin{enumerate}
\item $A$ is Ramsey;
\item its indices have a partition $I\sqcup J$, with $|I|,|J|\ge3$,
such that each group is projectively algebraic.
\end{enumerate}
The implication \textup{(2)}$\Rightarrow$\textup{(1)} uses the sufficient
direction of OpenAI's Theorem~1.1~\cite{openai2026}; it is also a circle
case of Strinz's two-group theorem~\cite[Theorem C2]{strinz2026spherical}.
The implication \textup{(1)}$\Rightarrow$\textup{(2)} is proved here
without either direction of OpenAI's theorem.
\end{maintheorem}

\begin{maincorollary}\label{cor:seven}
For seven distinct circle points with $\trdeg_\Q M(A)=3$, the set is
Ramsey if and only if some four-point cross ratio is algebraic over $\Q$.
The sufficient implication uses OpenAI's Theorem~1.1, sufficient
direction; the necessary implication uses neither direction.
\end{maincorollary}

In a partition of seven indices into groups of size at least three the
sizes are $3$ and $4$. For a four-point group all orderings of its cross
ratio are fractional linear transforms of one of them. Thus the
equivalence in Corollary~\ref{cor:seven} follows immediately from
Theorem~\ref{thm:main}. Its sufficient implication holds without a
moduli-degree hypothesis, as already follows from Strinz's results.

\begin{maincorollary}\label{cor:palvolgyi}
Let $x,y,z$ be algebraically independent real numbers over $\Q$. The
seven circle points with finite parameters
\begin{equation}\label{eq:appendix-family}
 -2,\quad 0,\quad 2,\quad x,\quad y,\quad -4-x-2y,\quad z
\end{equation}
are non-Ramsey. This conclusion uses neither direction of OpenAI's
Theorem~1.1.
\end{maincorollary}

This is the family discussed after Theorem~A.8 in Appendix~E of
P\'alv\"olgyi~\cite{palvolgyi2026}. His parameters are
$s_1=x$, $s_2=y$, $s_3=-4-s_1-2s_2$, and $s_4=z$. That discussion
leaves the Ramsey status open and imposes a neighbourhood condition for
the failure of its derivation-certificate method. The corollary above
has no neighbourhood condition. Section~\ref{sec:appendix-example}
gives a short cross-ratio proof and an independent finite second-order
certificate.

\begin{maintheorem}\label{thm:elliptic}
For each transcendental real $X>\sqrt[3]{2}$ and either choice
$Y=\pm\sqrt{X^3-2}$, the recurrence in
\eqref{eq:elliptic-map} produces seven distinct circle parameters
$X_0,\ldots,X_6$. Their circle set is Ramsey, its moduli field has
transcendence degree one, and all its four-point cross ratios are
transcendental. The Ramsey assertion uses the sufficient direction of
OpenAI's Theorem~1.1.
\end{maintheorem}

This example shows that the degree-three hypothesis in the necessary
implication of Corollary~\ref{cor:seven} cannot simply be omitted.
The existence of Ramsey seven-point sets with no algebraic cross ratio
is already exhibited by Strinz~\cite[Section~7, ``Non-two-group Ramsey
sets'']{strinz2026spherical}, using different and simpler rational
families, including
\[
 \{0,\tfrac12,\tfrac13,t,t+1,2t,1/t\}
 \qquad(t\text{ real and transcendental}).
\]
That positive conclusion also uses the sufficient direction of
OpenAI's Theorem~1.1.
The elliptic example is included as a concrete boundary construction.

\subsection*{Public results used for context}
OpenAI's \emph{A classification of finite Euclidean Ramsey
configurations}, dated September~23, 2026, states a necessary and
sufficient tensor condition in Theorem~1.1~\cite{openai2026}.
We use the release at commit
\texttt{adc7f1241b42e322a6451854ab7e4b4c146bf78a}, result~172.
The corresponding Lean development includes
\ancfile{fieldCriterion_ramsey} and \ancfile{classification}
in \ancfile{OAI/Combinatorics/EuclideanRamsey/Main.lean}
at that commit~\cite{openai2026lean}. Its source supplies both
directions; no Lean kernel replay is part of the present note.

Jones's October~7, 2026 note gives an explicit tensor certificate for
any six concyclic points~\cite[Section~2]{jones2026}. Composed with
the sufficient direction of OpenAI's Theorem~1.1, this proves that every circle set with
at most six points is Ramsey. Strinz's October~15, 2026 release states
this conclusion as Theorem~B(a) and supplies a Lean development
composing the six-point construction with the sufficient direction of
OpenAI's Theorem~1.1
\cite{strinz2026spherical}. Thus the six-point result is background
here.

P\'alv\"olgyi's September~20, 2026 preprint states in Theorem~A.3 that,
for any fixed distinct real $t_1,t_2,t_3$, if $t_4,\ldots,t_7$ are
algebraically independent over $\Q(t_1,t_2,t_3)$, the seven circle
points are non-Ramsey~\cite{palvolgyi2026}. That negative result has a
derivation and colouring proof and uses neither direction of OpenAI's
Theorem~1.1. Its applicability is stronger than independence of
all seven absolute parameters, but does not cover
\eqref{eq:appendix-family}.

Strinz's M\"obius-invariance theorem, released October~13, 2026,
states that a finite spherical set is Ramsey if and only if its image
under a M\"obius transformation of its containing sphere is Ramsey
\cite[Theorem C]{strinz2026mobius}. Its published argument uses both
directions of OpenAI's Theorem~1.1.
The October~15 spherical note gives two further sufficient conditions:
two M\"obius images of finite algebraic subsets of a sphere have Ramsey
union (Theorem~C2), and an arbitrary finite algebraic circle set together
with at most three arbitrary circle points is Ramsey (Theorem~B(b))
\cite{strinz2026spherical}. Both positive conclusions use the sufficient
direction of OpenAI's Theorem~1.1. Their circle case explains the sufficient halves
of Theorem~\ref{thm:main} and Corollary~\ref{cor:seven}. The explicit
certificate is recalled in Section~\ref{sec:positive}.

For comparison, on $S^{d-1}$ with $d\ge3$, algebraically independent
stereographic parameters give the generic threshold
$\binom{d+2}{2}$: fewer points are Ramsey and that many or more are
non-Ramsey. Strinz's ledger SR-C3 and SR-C29 records the positive half
as a consequence of OpenAI's Proposition~7.3, using the sufficient
direction of its Theorem~1.1; the negative boundary is P\'alv\"olgyi's Theorem~A.12 and
uses neither direction~\cite{strinz2026spherical,palvolgyi2026}.
On the circle, the relevant generic negative size is seven, and the
classification below concerns the degree-three locus within its moduli.

The proofs occupy Sections~\ref{sec:tensors}--\ref{sec:classification}.
Section~\ref{sec:appendix-example} treats Corollary~\ref{cor:palvolgyi},
and Section~\ref{sec:elliptic} proves Theorem~\ref{thm:elliptic}.
The ancillary files record exact finite checks; the structural
necessity argument itself is given over the full moduli field.
